\bookchapter{Capítulo 7: Análisis completo de las funciones avanzadas}

\chapterintro{La tecnología avanza sin cesar y las funciones del teléfono evolucionan continuamente. Comunicación por satélite, comunicación sin red, detección de accidentes de tráfico, tecnología SparkLink... estas funciones que suenan a ciencia ficción ya forman parte de nuestra vida. Conozcamos estas tecnologías punteras.}

\booksection{I. Comunicación por satélite: pedir auxilio sin señal}

La comunicación por satélite es uno de los grandes avances recientes en el ámbito de los teléfonos. Cuando estás en plena naturaleza, en lo más profundo de la montaña y el teléfono no tiene señal, la comunicación por satélite puede ser un salvavidas.

\booksubsection{Principio de funcionamiento:}

El teléfono establece conexión directa con satélites de comunicaciones en órbita y envía mensajes de auxilio. Actualmente admite principalmente el envío de SMS e información de ubicación; no permite llamadas.

\booksubsection{Modelos compatibles:}

\begin{itemize}
  \item \textbf{Apple}: A partir de la serie iPhone 14, se soporta el envío de SOS de emergencia por satélite. La serie iPhone 17 amplía el soporte para enviar iMessage y SMS por satélite, con cobertura en más regiones.
  \item \textbf{Huawei}: desde la serie Mate 50 admite mensajes por satélite BeiDou (solo texto); la serie Mate 60 Pro añade llamadas por satélite Tiantong (se pueden hacer llamadas); la serie Mate 70 refuerza aún más la capacidad de comunicación por satélite.
  \item \textbf{Xiaomi/OPPO/vivo}: desde 2024, los modelos de gama alta van incorporando progresivamente la función de mensajes por satélite BeiDou.
  \item \textbf{Tendencia}: la comunicación por satélite está pasando de ser una función exclusiva de la gama alta a popularizarse; en el futuro, los modelos de gama media también podrían admitirla. El sistema de satélites chino BeiDou es el más ampliamente compatible con los teléfonos nacionales.
\end{itemize}

\booksubsection{Escenarios de uso:}

\begin{itemize}
  \item Aventuras al aire libre, accidentes de montañismo.
  \item Emergencias en el mar.
  \item Interrupción de las comunicaciones terrestres por desastres naturales.
\end{itemize}

\booksubsection{Modo de empleo:}

\begin{enumerate}
  \item Asegúrate de estar en zona abierta, con vista al cielo.
  \item Activa la función de comunicación por satélite del teléfono (suele estar en la interfaz de llamada de emergencia).
  \item Sigue las indicaciones de la pantalla y orienta el teléfono hacia el satélite.
  \item Envía el mensaje de auxilio predefinido y la ubicación.
\end{enumerate}

\begin{tcbinfo}{Sabías que}
La comunicación por satélite requiere una antena de hardware especial en el teléfono; no todos los teléfonos la admiten. Además, suele requerir una suscripción de pago, aunque la función de emergencia suele ser gratuita.

\end{tcbinfo}

\booksection{II. Comunicación sin red: una nueva forma de comunicación a corta distancia}

La “comunicación sin red” se refiere a la tecnología que permite que dos teléfonos se comúniquen directamente sin red móvil ni Wi-Fi.

\booksubsection{Principio técnico:}

Usa Bluetooth o frecuencias de radio dedicadas para establecer una conexión punto a punto a corta distancia (normalmente de cientos de metros a pocos kilómetros).

\booksubsection{Escenarios de aplicación:}

\begin{itemize}
  \item \textbf{Actividades al aire libre}: comunicación entre miembros del grupo en senderismo o acampada.
  \item \textbf{Rescate en catástrofes}: cuando se interrumpen las comunicaciones tras un terremoto o una inundación.
  \item \textbf{Grandes eventos}: conciertos y eventos deportivos con mucha gente y mala señal.
\end{itemize}

\booksubsection{Plataformas compatibles:}

\begin{itemize}
  \item \textbf{Apple}: "NameDrop" y "Encontrar amigos sin conexión" en iOS 17.
  \item \textbf{Huawei}: la App Canghai admite llamadas sin red.
  \item \textbf{Xiaomi}: algunas funciones del Xiaomi Miaoxiang Center.
\end{itemize}

\booksection{III. Detección de accidentes de tráfico: salvar vidas en el momento crítico}

La detección de accidentes usa los sensores del teléfono para detectar accidentes graves y llamar automáticamente al teléfono de emergencia.

\booksubsection{Principio de funcionamiento:}

El teléfono combina varios sensores —acelerómetro, giroscopio, barómetro, micrófono— para determinar si se ha producido una colisión grave. Al detectar un accidente, el teléfono:

\begin{enumerate}
  \item Muestra una cuenta atrás y pregunta si estás bien.
  \item Si no respondes cuando termina la cuenta atrás, llama automáticamente al teléfono de emergencia.
  \item Envía tu ubicación a los contactos de emergencia.
\end{enumerate}

\booksubsection{Modelos compatibles:}

\begin{itemize}
  \item \textbf{Apple}: iPhone 14 y posteriores.
  \item \textbf{Android}: algúnos modelos de gama alta lo van incorporando.
\end{itemize}

\booksubsection{Precauciones:}

\begin{itemize}
  \item La función debe estar activada para usarse.
  \item Puede dispararse por error (por ejemplo, en montañas rusas o deportes intensos); basta con cancelarlo a tiempo.
  \item Depende de sensores y algoritmos; no es 100 \% fiable, pero añade una capa más de protección.
\end{itemize}

\booksection{IV. Tecnología SparkLink: la conexión inalámbrica de nueva generación}

SparkLink (Xingling) es una tecnología de comunicación inalámbrica de nueva generación impulsada por Huawei, con el objetivo de sustituir a Bluetooth y Wi-Fi en determinados escenarios. Supera a las soluciiones actuales en indicadores clave: la velocidad de transmisión puede alcanzar 900 Mbps (muy por encima de los 50 Mbps de Bluetooth), la latencia es de solo 20 $\mu$ s (Bluetooth, unos 15 ms) y una sola red admite 4096 dispositivos conectados simultáneamente (Bluetooth, normalmente menos de 8).

\booksubsection{Características técnicas:}

\begin{itemize}
  \item \textbf{Latencia ultrabaja}: latencia de extremo a extremo de solo 20 microsegundos (0,02 ms), dos órdenes de magnitud inferior a Bluetooth; la sincronización de juego y música es prácticamente imperceptible.
  \item \textbf{Alta velocidad}: velocidad máxima de un solo enlace de 900 Mbps, decenas de veces superior a la del audio Bluetooth; permite transmitir audio sin pérdida e incluso vídeo en alta definición.
  \item \textbf{Bajo consumo}: el consumo en reposo es similar al de Bluetooth, pero la eficiencia de transmisión es mayor, por lo que el consumo por unidad de datos es menor.
  \item \textbf{Amplia cobertura + alta concurrencia}: distancia de conexión más de un 50 \% mayor que Bluetooth; una sola red admite miles de dispositivos conectados a la vez.
\end{itemize}

\booksubsection{Escenarios de aplicación:}

\begin{itemize}
  \item \textbf{Auriculares inalámbricos}: menor latencia, sincronización perfecta entre imagen y sonido al jugar.
  \item \textbf{Smartwatch}: conexión más estable.
  \item \textbf{Hogar inteligente}: más dispositivos conectados a la vez.
  \item \textbf{Interconexión con el coche}: conexión más fluida entre teléfono y sistema del coche.
\end{itemize}

\booksubsection{Dispositivos compatibles:}

Actualmente, principalmente dispositivos del ecosistema Huawei, como la serie Huawei Mate 60 y algúnos auriculares y relojes Huawei. En el futuro, más marcas lo admitirán.

\textbf{Perspectiva tecnológica}\par

SparkLink representa el futuro de la comunicación inalámbrica. A medida que se sumen más fabricantes, veremos más dispositivos compatibles y disfrutaremos de una mejor experiencia inalámbrica.

\begin{tcbpractice}{Comprueba si tu teléfono admite estas funciones}
\begin{enumerate}
  \item Busca "satélite" en la configuración para ver si hay opciones relaciónadas con la comunicación por satélite.
  \item Busca "detección de accidentes" o "SOS de emergencia" y comprueba si están activadas.
  \item Si tu teléfono es Huawei, comprueba si admite la conexión de dispositivos SparkLink.
  \item Conoce cómo se usan estas funciones; en un momento crítico pueden ser de gran ayuda.
\end{enumerate}

\end{tcbpractice}
